% Part: propositional-logic
% Chapter: syntax-and-semantics
% Section: valuations-sat

\documentclass[../../../include/open-logic-section]{subfiles}

\begin{document}

\olfileid{pl}{syn}{val}

\olsection{\usetoken{P}{valuation} आणि पूर्ति}

\begin{defn}[!!^{valuation}s]
``सत्य'' आणि ``असत्य'' ही दोन सत्यतामूल्ये असलेला संच
$\{\True, \False\}$ असू द्या. $\Lang{L_0}$ साठीचे
\emph{!!{valuation}} म्हणजे भाषेतील !!{propositional variable}s ना
$\True$ किंवा $\False$ यांपैकी एक मूल्य देणारे फलन~$\pAssign{v}$;
म्हणजेच $\pAssign{v} \colon \PVar \to \{\True, \False \}$.
\end{defn}

\begin{defn}
  दिलेले !!a{valuation}~$\pAssign{v}$ असताना, मूल्यन फलन
  $\pValue{v} \colon \Frm[L_0] \to \{\True, \False \}$ पुढीलप्रमाणे
  विगमनाने परिभाषित करू:
  \begin{align*}
    \iftag{prvFalse}{\pValue{v}(\lfalse) & = \False; \\}{}
    \iftag{prvTrue}{\pValue{v}(\ltrue) & = \True; \\}{}
    \pValue{v}(\Obj p_n) & = \pAssign{v}(\Obj p_n); \\
    \iftag{prvNot}{\pValue{v}(\lnot !A) & = \begin{cases}
        \True & \text{जर } \pValue{v}(!A) = \False;\\
        \False & \text{अन्यथा.}
      \end{cases} \\ }{}
    \iftag{prvAnd}{\pValue{v}(!A \land !B) & = \begin{cases}
        \True &
        \text{जर $\pValue{v}(!A) = \True$ आणि $\pValue{v}(!B) = \True$;}\\
        \False &
        \text{जर $\pValue{v}(!A) = \False$ किंवा $\pValue{v}(!B) = \False$}.
      \end{cases}\\}{}
    \iftag{prvOr}{\pValue{v}(!A \lor !B) & = \begin{cases}
        \True &
        \text{जर $\pValue{v}(!A) = \True$ किंवा $\pValue{v}(!B) = \True$;}\\
        \False &
        \text{जर $\pValue{v}(!A) = \False$ आणि $\pValue{v}(!B) = \False$}.
      \end{cases}\\}{}
    \iftag{prvIf}{\pValue{v}(!A \lif !B) & = \begin{cases}
        \True &
        \text{जर $\pValue{v}(!A) = \False$ किंवा $\pValue{v}(!B) = \True$;}\\
        \False &
        \text{जर $\pValue{v}(!A) = \True$ आणि $\pValue{v}(!B) = \False$}.
      \end{cases}\\}{}
    \iftag{prvIff}{\pValue{v}(!A \liff !B) & = \begin{cases}
        \True &
        \text{जर $\pValue{v}(!A) = \pValue{v}(!B)$;}\\
        \False &
        \text{जर $\pValue{v}(!A) \neq \pValue{v}(!B)$}.
      \end{cases}}{}
\end{align*}
\end{defn}

\begin{explain}
या बाबींना पुढील सत्यताकोष्टके अनुरूप आहेत:
\begin{center}
\iftag{prvNot}{
  \begin{tabular}{|c||c|} \hline
    $!A$ & $ \lnot !A$ \\
    \hline \hline
    $\True$ & $\False$ \\
    $\False$ & $\True$ \\
    \hline
  \end{tabular}
}{}
\iftag{prvAnd}{
  \begin{tabular}{|cc||c|} \hline
    $!A$ & $!B$ & $!A \land !B$ \\
    \hline \hline
    $\True$ & $\True$ & $\True$ \\
    $\True$ & $\False$ & $\False$ \\
    $\False$ & $\True$ & $\False$ \\
    $\False$ & $\False$ & $\False$ \\
    \hline
  \end{tabular}
}{}
\iftag{prvOr}{
  \begin{tabular}{|cc||c|} \hline
    $!A$ & $!B$ & $!A \lor !B$ \\
    \hline \hline
    $\True$ & $\True$ & $\True$ \\
    $\True$ & $\False$ & $\True$ \\
    $\False$ & $\True$ & $\True$ \\
    $\False$ & $\False$ & $\False$ \\
    \hline
  \end{tabular}
}{}%
\iftag{notprvNot,notprvAnd,notprvOr,notprvIf}{}{\\[1em]}
\iftag{prvIf}{
  \begin{tabular}{|cc||c|} \hline
    $!A$ & $!B$ & $!A \lif !B$ \\
    \hline \hline
    $\True$ & $\True$ & $\True$ \\
    $\True$ & $\False$ & $\False$ \\
    $\False$ & $\True$ & $\True$ \\
    $\False$ & $\False$ & $\True$ \\
    \hline
  \end{tabular}
}{}
\iftag{prvIff}{
  \begin{tabular}{|cc||c|} \hline
    $!A$ & $!B$ & $!A \liff !B$ \\
    \hline \hline
    $\True$ & $\True$ & $\True$ \\
    $\True$ & $\False$ & $\False$ \\
    $\False$ & $\True$ & $\False$ \\
    $\False$ & $\False$ & $\True$ \\
    \hline
  \end{tabular}
}{}
\end{center}
\end{explain}

\begin{prob}
$\Lang{L_0}$ मध्ये $\diamondsuit$ हा त्रिस्थानी संयोजक वाढवण्याचा विचार
करा; त्याचे मूल्यन पुढीलप्रमाणे दिले आहे:
\begin{gather*}
  \pValue{v}(\diamondsuit( !A , !B , !C ) ) = \begin{cases}
    \pValue{v}( !B ) &
    \text{जर $\pValue{v}(!A) = \True$;}\\
    \pValue{v}( !C ) &
    \text{जर $\pValue{v}(!A) = \False $}.
  \end{cases}
\end{gather*}
या संयोजकाचे सत्यताकोष्टक लिहा.
\end{prob}

\begin{thm}[स्थानिक निर्धारण]
\ollabel{thm:LocalDetermination} $\pAssign{v_1}$ आणि
$\pAssign{v_2}$ ही अशी !!{valuation}s आहेत की $!A$ मध्ये आस्थित असलेल्या
!!{propositional
variable}s ना ती समान मूल्ये देतात; म्हणजे, एखादे $\Obj p_n$ हे
!!{formula}~$!A$ मध्ये आस्थित असेल तेव्हा $\pAssign{v_1}(\Obj p_n) =
\pAssign{v_2}(\Obj p_n)$. मग $\pValue{v_1}$ आणि $\pValue{v_2}$ हीही
$!A$ वर एकमत असतात; म्हणजे $\pValue{v_1}(!A) = \pValue{v_2}(!A)$.
\end{thm}

\begin{proof}
$!A$ वरील विगमनाने.
\end{proof}

\begin{defn}[पूर्ति]
\ollabel{defn:satisfaction} आपण
  \emph{!!a{valuation}~$\pAssign{v}$ कडून !!a{formula}~$!A$ ची पूर्ति},
  $\pSat{v}{!A}$, ही संकल्पना पुढीलप्रमाणे विगमनाने परिभाषित करू शकतो.
  (``$\pSat{v}{!A}$ नाही'' हे दर्शवण्यासाठी आपण $\pSat/{v}{!A}$ लिहितो.)
\begin{enumerate}
\tagitem{prvFalse}{%
  \indcase{!A}{\lfalse}{$\pSat/{v}{\indfrm}$.}}{}

\tagitem{prvTrue}{%
  \indcase{!A}{\ltrue}{$\pSat{v}{\indfrm}$.}}{}

\item \indcase{!A}{\Obj p_i}{$\pSat{v}{\indfrm}$ तेव्हा आणि केवळ
  तेव्हाच, जेव्हा $\pAssign{v}(\Obj p_i) = \True$.}

\tagitem{prvNot}{%
  \indcase{!A}{\lnot !B}{$\pSat{v}{\indfrm}$ तेव्हा आणि केवळ
    तेव्हाच, जेव्हा $\pSat/{v}{!B}$.}}{}

\tagitem{prvAnd}{%
  \indcase{!A}{(!B \land !C)}{$\pSat{v}{\indfrm}$ तेव्हा आणि केवळ
    तेव्हाच, जेव्हा $\pSat{v}{!B}$ आणि $\pSat{v}{!C}$.}}{}

\tagitem{prvOr}{%
  \indcase{!A}{(!B \lor !C)}{$\pSat{v}{\indfrm}$ तेव्हा आणि केवळ
    तेव्हाच, जेव्हा $\pSat{v}{!B}$ किंवा $\pSat{v}{!C}$ (किंवा दोन्ही).}}{}

\tagitem{prvIf}{%
  \indcase{!A}{(!B \lif !C)}{$\pSat{v}{\indfrm}$ तेव्हा आणि केवळ
    तेव्हाच, जेव्हा $\pSat/{v}{!B}$ किंवा $\pSat{v}{!C}$ (किंवा दोन्ही).}}{}

\tagitem{prvIff}{%
  \indcase{!A}{(!B \liff !C)}{$\pSat{v}{\indfrm}$ तेव्हा आणि केवळ
    तेव्हाच, जेव्हा $\pSat{v}{!B}$ आणि $\pSat{v}{!C}$ ही दोन्ही, किंवा
    $\pSat{v}{!B}$ आणि $\pSat{v}{!C}$ यांपैकी कोणतेही नाही.}}{}
\end{enumerate}
$\Gamma$ हा !!{formula}s चा संच असेल, तर $\pSat{v}{\Gamma}$ तेव्हा आणि
केवळ तेव्हाच, जेव्हा प्रत्येक~$!A \in \Gamma$ साठी $\pSat{v}{!A}$.
\end{defn}

\begin{prop}\ollabel{prop:sat-value}
  $\pSat{v}{!A}$ तेव्हा आणि केवळ तेव्हाच, जेव्हा
  $\pValue{v}(!A) = \True$.
\end{prop}

\begin{proof}
  $!A$ वरील विगमनाने.
\end{proof}

\begin{prob}
  \olref[pl][syn][val]{prop:sat-value} सिद्ध करा.
\end{prob}

\end{document}
